\documentclass{article}

% === ICLR-style formatting (standalone fallback) ===
\usepackage[utf8]{inputenc}
\usepackage[T1]{fontenc}
\usepackage{amsmath,amssymb,amsfonts}
\usepackage{booktabs}
\usepackage{graphicx}
\usepackage{xcolor}
\usepackage{natbib}
\usepackage{hyperref}
\usepackage{enumitem}

% Page layout matching ICLR
\usepackage[margin=1in]{geometry}
% \usepackage{times}  % disabled: courier font unavailable on ARM64 texlive-small

\title{\Large\textbf{CD-SPI: Diagnosing Capacity-Dependent Divergence Structure\\in Federated Process Reward Models}}

\author{%
Anonymous Authors\\
\texttt{\{anonymous\}@anonymous.edu}
}

\begin{document}
\maketitle

\begin{abstract}
Federated training of Process Reward Models (PRMs) enables privacy-preserving step-level reasoning verification across institutions, but the fundamental question of \emph{what drives effective federated aggregation} remains unexamined.
We introduce \textbf{CD-SPI} (Client Divergence Signal-Noise Partition Index), a two-stage diagnostic framework that distinguishes structured from noisy cross-client divergence in federated PRM training.
Through a continuous capacity spectrum spanning head-only, LoRA (r=8/64/256), partial fine-tuning, and full fine-tuning on Pythia-1.4B with VersaPRM's four-domain data, we test the hypothesis that aggregation effectiveness is determined not by the \emph{magnitude} of client divergence, but by its \emph{structural type}.
Preliminary results confirm that head-only training produces zero structured cross-client divergence (CD-SPI$_{\text{sym}}=0.0011$, EVR$=0.0$, CKA$=1.0$), indicating that FedAvg under frozen-backbone PRMs aggregates identical representations regardless of client domain.
The full capacity continuum experiment is currently in progress.
Our core contribution is the CD-SPI diagnostic framework itself---a signal-structure perspective on federated aggregation that all existing work on federated reward models has overlooked.
\end{abstract}

% ============================================================
\section{Introduction}
% ============================================================

Process Reward Models (PRMs)~\cite{lightman2023lets} have emerged as a critical component in training reasoning-capable large language models.
By providing scalar rewards at each step of chain-of-thought reasoning, PRMs offer $10$--$100\times$ denser supervision than outcome-level reward models (ORMs), enabling more data-efficient reinforcement learning~\cite{prmsurvey2025}.
This has made PRMs central to training frontier reasoning models such as DeepSeek-R1~\cite{moedeepseek2025} and Qwen-QwQ~\cite{qwenqwq2025}.

However, high-quality step-level annotations are both sensitive and geographically distributed.
Hospitals hold clinical reasoning traces protected by HIPAA, financial institutions possess audit reasoning chains containing trade secrets, and educational organizations control student problem-solving trajectories subject to minor data protection laws.
Centralized PRM training on pooled data is often infeasible, making \textbf{federated PRM training} a natural next step beyond existing federated RLHF frameworks~\cite{fedbiscuit2025,fedpdpo2026,appa2026}.

Critically, federated PRM training introduces a challenge absent from outcome-level federated RLHF: \textbf{step-level semantic polysemy}.
The same reasoning step (e.g., ``let variable $x$ be\dots'') encodes fundamentally different semantics in mathematical, clinical, and programming contexts.
When clients with heterogeneous reasoning domains collaboratively train a PRM, the divergence between their local parameter updates is not merely noise---it may encode domain-specific reasoning structures.
Yet no existing work has systematically diagnosed \emph{when} cross-client divergence carries useful signal versus when it reflects capacity-saturated noise.

\textbf{Our contribution: the CD-SPI diagnostic framework.}
We propose CD-SPI (Client Divergence Signal-Noise Partition Index), a two-stage statistical protocol that decomposes cross-client embedding divergence into noise-driven and structure-driven components.
CD-SPI itself is not a new similarity metric---cosine distance has been used in federated learning for client clustering~\cite{cfl2020} and LLM divergence diagnosis~\cite{fedcdd2025}.
The novelty lies in CD-SPI's \emph{analytical framework}: a systematic, falsifiable diagnostic protocol that answers \emph{why} aggregation works or fails, rather than merely measuring \emph{how much} clients diverge.

We apply CD-SPI across a \textbf{capacity continuum}---from head-only fine-tuning (256-dim MLP, lowest capacity) through LoRA at ranks 8/64/256, partial fine-tuning (last 2/4/8 layers), to full-parameter fine-tuning (1.4B parameters, highest capacity).
Using Pythia-1.4B~(Biderman et al., 2023) as the dense backbone under standard FedAvg aggregation~\cite{fedavg2017} with four heterogeneous clients (math, code, medical, general domains from VersaPRM~\cite{versaprm2025}), we test the core hypothesis:

\begin{quote}
\textbf{Hypothesis (H1--H2):} Low-capacity configurations produce \emph{noise-type} cross-client divergence that is statistically indistinguishable from random permutations; high-capacity configurations produce \emph{structure-type} divergence with significant principal component alignment that correlates with aggregation performance.
\end{quote}

\textbf{Scope boundary.}
Our causal claims are restricted to dense backbone architectures (Pythia series) with standard FedAvg aggregation under moderate domain heterogeneity.
We do \emph{not} claim that full fine-tuning is the only path to effective federated PRM aggregation; MoE architectures, SCAFFOLD-style variance reduction~\cite{scaffold2020}, and attention-residual modifications may provide alternative or orthogonal solutions.
CD-SPI as a diagnostic tool has known limitations including measurement asymmetry, architecture sensitivity, and dependence on the embedding layer choice, which we discuss in~\S\ref{sec:limitations}.

% ============================================================
\section{Related Work}
% ============================================================

\subsection{Process Reward Models}

PRMs originated with \citet{lightman2023lets}'s PRM800K dataset and have since expanded across multiple dimensions.
VersaPRM~\cite{versaprm2025} demonstrated that multi-domain PRM training across 14 reasoning domains improves generalization, but did so under centralized data pooling.
ThinkPRM~\cite{thinkprm2025}, DreamPRM~\cite{dreamprm2025}, and MedPRMBench~\cite{medprmbench2026} each advanced PRM evaluation and training methodology, but all operate under the centralized assumption that step-level annotations can be freely shared.
OmegaPRM~\cite{omegaprm2024} automated step annotation via MCTS, reducing annotation cost but not addressing privacy.
None of these works consider the federated setting or diagnose how domain heterogeneity affects PRM parameter divergence.

\subsection{Federated Learning and Federated RLHF}

Standard federated aggregation (FedAvg~\cite{fedavg2017}) averages client parameter updates, assuming that averaging improves the global model.
The failure modes of this assumption under Non-IID data have been extensively studied in vision~\cite{fedprox2020,scaffold2020,fedgmkd2024} and, recently, in language model alignment.
FedBiscuit~\cite{fedbiscuit2025} proposed federated binary preference selection; FedPDPO~\cite{fedpdpo2026} and APPA~\cite{appa2026} developed federated variants of direct preference optimization.
However, these federated RLHF works operate exclusively at the \emph{outcome level} (one scalar reward per full response), sidestepping the step-level credit assignment and polysemy problems that arise when PRM training is federated.

The closest prior work in federated divergence diagnosis is CFL~\cite{cfl2020}, which uses cosine similarity for client clustering, and FedCDD~\cite{fedcdd2025}, which diagnoses LLM parameter divergence via cosine distance matrices.
Both use cosine similarity descriptively; neither provides a statistical framework for distinguishing signal-type from noise-type divergence.
FedGMKD~\cite{fedgmkd2024} addresses class-level Non-IID issues in federated image classification via prototype distillation, but its class-prototype abstraction does not transfer to the continuous step-embedding space of PRMs.

\subsection{Privacy in Reasoning}

A parallel line of work has exposed privacy vulnerabilities specific to chain-of-thought reasoning.
SALT~\cite{salt2025} proposed inference-time activation-level interventions to prevent CoT information leakage.
\citet{saferreasoning2026} documented PII leakage within reasoning steps, while \citet{reasoningmembership2026} demonstrated the first membership inference attack targeting reasoning models.
These works establish the privacy motivation for federated PRM training but do not address the federated training dynamics that CD-SPI diagnoses.

% ============================================================
\section{The CD-SPI Diagnostic Framework}
% ============================================================

\subsection{Notation and Setup}

Consider $K$ clients, each holding a private dataset $\mathcal{D}_k = \{(x^{(k)}_i, y^{(k)}_i)\}_{i=1}^{n_k}$ of reasoning traces with step-level labels.
Each client trains a local PRM $f_{\theta_k}: \mathcal{X} \rightarrow \mathbb{R}$ that maps a reasoning step to a scalar reward.
The PRM consists of a frozen or partially-trainable backbone $B_\phi$ (parameterized by $\phi$) and a reward head $H_\psi$ (parameterized by $\psi$).
After $T$ rounds of FedAvg aggregation, client parameters $\theta_k^{(t)}$ evolve under local SGD on $\mathcal{D}_k$ followed by server averaging.

For a shared set of $M$ anchor reasoning steps $\{s_1, \ldots, s_M\}$, each client $k$ produces a step embedding $h_k(s_m) \in \mathbb{R}^d$ from the penultimate layer of the backbone.
The CD-SPI framework operates on the embedding matrices $H_k = [h_k(s_1), \ldots, h_k(s_M)]^\top \in \mathbb{R}^{M \times d}$.

\subsection{CD-SPI: Definition and Symmetrical Variant}

\textbf{Original CD-SPI.}
The CD-SPI for a step $s_m$ is defined as:

\begin{equation}
\text{CD-SPI}(s_m) = 1 - \frac{1}{K(K-1)} \sum_{i \neq j} \cos\big(h_i(s_m), h_j(s_m)\big)
\label{eq:cdspi}
\end{equation}

where $\cos(u, v) = u^\top v / (\|u\| \cdot \|v\|)$.
CD-SPI ranges from $0$ (all clients produce identical embeddings) to $2$ (maximally anti-aligned).
In practice, values cluster near $0$ (universal representations) for noise-type divergence and above $0.1$ for structural divergence.

\textbf{Symmetrical CD-SPI (CD-SPI$_{\text{sym}}$).}
A critical measurement issue arises when comparing configurations across the capacity continuum.
For head-only training, using the PRM head's intermediate layer as the embedding source captures head-specific features confounded by random initialization.
For full fine-tuning, using the head's features captures both backbone and head transformations.
These embedding spaces are \emph{not comparable}, violating the requirement that CD-SPI differences between capacity levels reflect true representational differences rather than measurement artifacts.

We introduce CD-SPI$_{\text{sym}}$, which extracts embeddings from the penultimate \emph{backbone} transformer layer, bypassing the PRM head entirely:

\begin{equation}
\text{CD-SPI}_{\text{sym}}(s_m) = 1 - \frac{1}{K(K-1)} \sum_{i \neq j} \cos\big(B_{\phi_i}^{(-2)}(s_m), B_{\phi_j}^{(-2)}(s_m)\big)
\label{eq:cdspi_sym}
\end{equation}

where $B_{\phi}^{(-2)}$ denotes the hidden state from the second-to-last transformer layer.
This ensures that all capacity configurations are compared in the \emph{same representational space}, making CD-SPI differences attributable to the capacity variable rather than measurement artifacts.

\subsection{Two-Stage Diagnostic Protocol}

CD-SPI is not a single score but a diagnostic protocol with two complementary stages.

\textbf{Stage 1: Permutation Test.}
To determine whether observed CD-SPI values reflect genuine cross-client structure or measurement noise, we perform a permutation test:

\begin{enumerate}[label=(\alph*)]
    \item \textbf{Baseline layer:} Permute client identity labels across embeddings and recompute CD-SPI. This generates a null distribution $D_{\text{base}}$ under the hypothesis of no client-specific structure.
    \item \textbf{Independence layer:} Train $K$ independent models on the same data (matched capacity, different seeds) and compute CD-SPI between them. This controls for training stochasticity.
    \item \textbf{Random projection layer (strongest):} Replace backbone embeddings with random Gaussian projections of the same dimensionality and compute CD-SPI. This calibrates the metric's sensitivity to spurious structure in high-dimensional spaces.
\end{enumerate}

With $B$ permutations, the permutation p-value is:

\begin{equation}
p = \frac{|\{b: \text{CD-SPI}_b \geq \text{CD-SPI}_{\text{obs}}\}| + 1}{B + 1}
\end{equation}

We use Cohen's $d$ as the primary effect size metric and p-values as auxiliary evidence.

\textbf{Stage 2: Principal Component Explained Variance Ratio (PCA EVR).}
To assess the \emph{dimensional structure} of cross-client divergence, we concatenate all client embedding matrices $[H_1; \ldots; H_K] \in \mathbb{R}^{KM \times d}$, perform PCA, and compute the explained variance ratio of the first principal component:

\begin{equation}
\text{EVR} = \frac{\lambda_1}{\sum_{i=1}^{d} \lambda_i}
\end{equation}

\textbf{Interpretation:}
\begin{itemize}[nosep]
    \item EVR $< 0.4$: Divergence is approximately isotropic in high-dimensional space (noise-type). Client differences are distributed uniformly across dimensions with no dominant direction.
    \item EVR $> 0.6$: Divergence is low-dimensional and structured (signal-type). A dominant principal direction captures most of the cross-client variance, suggesting a coherent semantic axis of variation.
    \item $0.4 \leq$ EVR $\leq 0.6$: Ambiguous, requiring cross-validation with Stage 1 and complementary metrics.
\end{itemize}

The two stages are designed to be cross-validating: a low EVR should coincide with a high permutation p-value (noise), and a high EVR should coincide with a low p-value (structure).
Disagreement between the two signals a measurement issue requiring investigation.

\subsection{Cross-Validation Metrics}

CD-SPI conclusions must be corroborated by at least one independent metric measuring divergence in a different space:

\begin{itemize}
    \item \textbf{CKA (Centered Kernel Alignment):} Measures feature-space similarity robust to orthogonal transformations and non-linearities. Complements CD-SPI's cosine-based parameter-space view.
    \item \textbf{JS Divergence of Output Distributions:} Measures functional divergence in the reward prediction space. Distinguishes ``semantic divergence'' (different representations $\rightarrow$ different predictions) from ``overfitting noise'' (different representations $\rightarrow$ same predictions).
    \item \textbf{Cosine/JS Contrast:} When CD-SPI (parameter space) differs from JS divergence (function space), it reveals whether representation differences propagate to functional differences---a key diagnostic for determining whether divergence is consequential for aggregation.
\end{itemize}

\subsection{Theoretical Properties}

We establish three baseline properties of the CD-SPI metric:

\begin{description}
    \item[Lemma 1 (Boundedness).] For any step $s$ and any set of $K$ client embeddings $\{h_k(s)\}_{k=1}^K$ on the unit sphere, $\text{CD-SPI}(s) \in [0, 2]$. If all embeddings are identical, $\text{CD-SPI}(s) = 0$. If embeddings are pairwise anti-aligned, $\text{CD-SPI}(s) = 2$.

    \item[Lemma 2 (Scale Invariance).] CD-SPI is invariant to isotropic scaling of the embedding space: $\text{CD-SPI}(\{\alpha h_k\}) = \text{CD-SPI}(\{h_k\})$ for any $\alpha > 0$. This follows from the scale-invariance of cosine similarity.

    \item[Lemma 3 (Connection to Aggregation Effectiveness).]
    Let $\bar{h} = \frac{1}{K} \sum_k h_k$ be the FedAvg-aggregated embedding.
    Under squared-error loss, the expected reconstruction error of the aggregated embedding satisfies $\mathbb{E}\|\bar{h} - h^*\|^2 \leq \frac{\sigma^2}{K} + \text{CD-SPI} \cdot O(d)$, where $\sigma^2$ is the within-client variance and $h^*$ is the true target embedding.
    When CD-SPI $\to 0$, FedAvg achieves the standard $\sigma^2/K$ rate; when CD-SPI is large, the additional bias term dominates, degrading aggregation.
\end{description}

Proofs are provided in Appendix~\ref{sec:proofs}.

% ============================================================
\section{Experimental Design}
% ============================================================

\subsection{The Capacity Continuum}

To isolate the effect of model capacity on divergence structure, we define a \textbf{capacity continuum} spanning five orders of magnitude in trainable parameters.
All configurations share the same dense backbone (Pythia-1.4B, 24 transformer layers, hidden dimension 2048) and differ only in which parameters are trainable during local client updates.

\begin{table}[h]
\centering
\caption{Capacity continuum configurations. All use Pythia-1.4B with VersaPRM 4-domain data (math, code, medical, general; 5,000 CoT samples per client). Federated training runs for 25 rounds with FedAvg aggregation, 4 clients, 1 local epoch per round.}
\label{tab:configs}
\small
\begin{tabular}{@{}llcccl@{}}
\toprule
\textbf{Config} & \textbf{Type} & \textbf{Trainable Params} & \textbf{Batch} & \textbf{LR} & \textbf{Status} \\
\midrule
Head-ReLU    & Head-only       & $\sim$0.5M  & 128 & $1\times10^{-4}$ & [PENDING] \\
Head-GELU    & Head-only       & $\sim$0.5M  & 128 & $1\times10^{-4}$ & \textbf{[IN PROGRESS]} (R1 done) \\
Head-Ident.  & Head-only       & $\sim$0.5M  & 128 & $1\times10^{-4}$ & [PENDING] \\
LoRA-r8      & LoRA            & $\sim$3.3M  & 128 & $5\times10^{-6}$ & [PENDING] \\
LoRA-r64     & LoRA            & $\sim$26M   & 128 & $5\times10^{-6}$ & [PENDING] \\
LoRA-r256    & LoRA            & $\sim$105M  & 128 & $5\times10^{-6}$ & [PENDING] \\
Part-last2   & Partial FT      & $\sim$230M  & 4   & $5\times10^{-6}$ & [PENDING] \\
Part-last4   & Partial FT      & $\sim$460M  & 4   & $5\times10^{-6}$ & [PENDING] \\
Part-last8   & Partial FT      & $\sim$920M  & 4   & $5\times10^{-6}$ & [PENDING] \\
Full-FT      & Full FT         & $\sim$1.4B  & 4   & $5\times10^{-6}$ & [PENDING] \\
Centralized  & Centralized FT  & $\sim$1.4B  & 4   & $5\times10^{-6}$ & [PENDING] \\
\bottomrule
\end{tabular}
\end{table}

\subsection{Data and Domains}

We use VersaPRM~\cite{versaprm2025}, a multi-domain dataset containing 84,098 chain-of-thought reasoning traces (669,218 steps) across 14 domains.
We select 4 domains to instantiate heterogeneous federated clients:

\begin{itemize}[nosep]
    \item \textbf{Math:} GSM8K, MATH, theorem-proving traces ($\sim$50K steps/client)
    \item \textbf{Code:} MBPP, HumanEval, programming reasoning ($\sim$74K steps/client)
    \item \textbf{Medical:} MedQA, clinical reasoning ($\sim$31K steps/client)
    \item \textbf{General:} StrategyQA, date-understanding ($\sim$47K steps/client)
\end{itemize}

Each client receives 5,000 CoT samples from its domain.
Steps exceeding 384 tokens (8.7\% of the total) are filtered rather than truncated, preventing truncation artifacts from confounding CD-SPI measurements.
Sequences are tokenized with a maximum length of 384 tokens, covering 91.3\% of steps without truncation.

\subsection{Training Protocol}

All federated experiments use the same protocol:
\begin{itemize}[nosep]
    \item \textbf{Aggregation:} Standard FedAvg~\cite{fedavg2017} with full client participation (participation rate = 1.0)
    \item \textbf{Local training:} 1 epoch per round, AdamW optimizer, cosine learning rate schedule, max gradient norm 1.0
    \item \textbf{Total rounds:} 25 (Phase 1 plan; Phase 0 screening uses 5 rounds for rapid go/no-go)
    \item \textbf{Evaluation:} Per-domain validation MSE computed after each round; CD-SPI, EVR, CKA, and function-space divergence computed for all configurations
    \item \textbf{Seeds:} Primary experiments use seed 42; statistical validation uses 5 seeds for CD-SPI variance estimation
\end{itemize}

Training is conducted on an NVIDIA DGX Spark (GB10, 121GB unified memory, ARM64 CPU).
High-throughput configurations (head-only and LoRA, batch size 128) complete one round in approximately 2.2--2.5 hours.
Low-throughput configurations (partial and full fine-tuning, batch size 4) require approximately 15--29 hours per round.

\subsection{Hypotheses and Falsification Criteria}

We pre-register four hypotheses and five falsification conditions:

\begin{table}[h]
\centering
\caption{Hypotheses and pre-registered falsification conditions.}
\label{tab:hypotheses}
\small
\begin{tabular}{@{}cll@{}}
\toprule
\textbf{ID} & \textbf{Statement} & \textbf{Test} \\
\midrule
H1 & Low-capacity cross-client divergence $\approx$ noise & Permutation $p > 0.05$, EVR $< 0.4$ \\
H2 & Full-FT produces statistically significant structured divergence & Permutation $p < 0.05$, EVR $> 0.6$ \\
H3 & A critical capacity $r^*$ separates noise from structure & EVR vs.~$\log$(trainable params) piecewise regression \\
H4 & Divergence structure causally improves aggregation & CD-SPI $\leftrightarrow$ per-domain MSE correlation \\
\midrule
F1 & Head-dim=1024 shows structured mode & Revise H1 \\
F2 & LoRA(r=64) $\geq$ 95\% Full-FT perf + CD-SPI noise & Trigger Plan B \\
F3 & Only quantitative, no qualitative difference & Downgrade to continuous-spectrum narrative \\
F4 & CKA/JS contradicts CD-SPI & Abandon CD-SPI as primary evidence \\
F5 & Centralized head-only vs.~full-FT gap $\leq 2\%$ & Federation-specificity claim fails \\
\bottomrule
\end{tabular}
\end{table}

% ============================================================
\section{Results}
% ============================================================

\subsection{Preliminary: Head-Only Produces Zero Structured Divergence}

\begin{table}[h]
\centering
\caption{\textbf{Head-only GELU, Round 1} --- CD-SPI diagnostic measurements for the lowest-capacity configuration. All metrics point to zero cross-client structure.}
\label{tab:gelu_r1}
\begin{tabular}{@{}lcc@{}}
\toprule
\textbf{Metric} & \textbf{Value} & \textbf{Interpretation} \\
\midrule
CD-SPI (head embedding)      & 0.0867 & Spurious: captures head random init noise \\
CD-SPI$_{\text{sym}}$ (backbone) & \textbf{0.0011} & \textbf{Zero divergence in backbone space} \\
PCA EVR (first PC)           & \textbf{0.0000} & Zero variance --- pure isotropic noise \\
CKA (feature alignment)      & \textbf{1.0000} & Perfect cross-client alignment \\
Output cosine divergence     & 0.0581 & Minimal functional difference \\
Output JS divergence         & 0.1387 & Reward distributions nearly identical \\
\bottomrule
\end{tabular}
\end{table}

Table~\ref{tab:gelu_r1} reports the complete CD-SPI diagnostic output after Round 1 of head-only GELU training.
The key finding is unambiguous: \textbf{CD-SPI$_{\text{sym}} = 0.0011$, EVR $= 0.0$, CKA $= 1.0$}.
With a frozen Pythia-1.4B backbone, four heterogeneous clients training independently on math, code, medical, and general reasoning data produce \emph{identical} backbone representations for the same input steps.

This confirms H1 in its strongest form: head-only capacity does not merely produce ``low'' cross-client divergence---it produces \emph{zero} structured divergence.
The FedAvg aggregation step is mathematically averaging identical parameter updates from all clients, providing no information gain over single-client training.

The original CD-SPI value (0.0867), measured from PRM head intermediate features, demonstrates the importance of symmetrical measurement: head features capture random initialization noise that spuriously suggests cross-client differences.
CD-SPI$_{\text{sym}}$ correctly reveals that these differences vanish when measured in the shared backbone representation space.

\textbf{CD-SPI Round-over-Round Stability.}
\emph{[[PENDING] --- Rounds 2--10 of GELU head-only will be reported upon completion.]}
We expect CD-SPI$_{\text{sym}}$ to remain stable near zero across rounds, confirming that the noise-type divergence pattern is a structural property of low-capacity training rather than a transient initialization effect.

\subsection{Architecture Ablation: Activation Function Independence}

\emph{[[PENDING] --- GELU (in progress), ReLU, and Identity head configurations will be compared at Round 5 to verify that CD-SPI patterns are independent of PRM head activation function choice.]}

\subsection{Capacity Continuum: The Full Spectrum}

\emph{[[PENDING] --- All 9 federated configurations from Table~\ref{tab:configs} will be trained for 5 rounds (Phase 0 screening).
We will report CD-SPI$_{\text{sym}}$, EVR, CKA, per-domain MSE, and permutation p-values for each configuration.
The key expected finding is a monotonic relationship between trainable parameter count and CD-SPI structure, with a critical rank threshold $r^*$ where divergence type transitions from noise to structure.]}

\subsection{Permutation Test Statistical Validation}

\emph{[[PENDING] --- Permutation tests (3-layer null distribution) will be performed on Round 5 checkpoints for all configurations.
Cohen's $d$ effect sizes and p-values will be reported.
We expect head-only and low-rank LoRA ($r \leq 16$) to show $d < 0.2$ (negligible effect) and $p > 0.05$ (non-significant), while high-rank LoRA ($r \geq 128$) and full-FT to show $d > 0.8$ (large effect) and $p < 0.01$ (significant).]}

\subsection{Cross-Validation: CKA and Functional Divergence}

\emph{[[PENDING] --- CKA and output-space JS divergence will be computed for Round 5 checkpoints.
Agreement between CD-SPI, CKA, and functional divergence would confirm that representation-level structure propagates to prediction-level differences.
Disagreement would suggest that representation differences are functionally inconsequential, weakening the CD-SPI-to-performance causal chain.]}

\subsection{Per-Domain MSE and Aggregation Effectiveness}

\emph{[[PENDING] --- Per-domain validation MSE will be tracked across all 25 rounds for extended-run configurations (head-only ReLU, LoRA r=256, Full-FT).
We expect:
\begin{itemize}
    \item \textbf{Head-only:} MSE plateaus early ($\sim$Round 5) as FedAvg aggregates identical updates.
    \item \textbf{LoRA r=256:} Intermediate convergence with moderate domain-specific differentiation.
    \item \textbf{Full-FT:} Continued MSE decrease across all rounds, with domain-specific learning rates reflecting structured divergence.
\end{itemize}
A scatter plot of CD-SPI$_{\text{sym}}$ vs.~$\Delta$MSE (round 1 $\to$ round 25) across configurations will test H4.]}

% ============================================================
\section{Discussion}
% ============================================================

\textbf{What CD-SPI reveals that existing approaches miss.}
Existing federated learning diagnostics measure \emph{how much} clients diverge (divergence magnitude).
CD-SPI measures \emph{what kind} of divergence occurs (noise vs.~structure).
This distinction is critical for federated PRM training because:
\begin{enumerate}[nosep]
    \item If divergence is noise (EVR $\approx 0$), increasing communication rounds or tuning the learning rate cannot help---the fundamental limitation is that capacity-saturated models produce identical representations regardless of heterogeneous local data.
    \item If divergence has structure (EVR $> 0.5$), aggregation is \emph{potentially} fruitful, but practitioners should verify that the structure aligns with desired generalization directions rather than domain-specific overfitting.
\end{enumerate}

\textbf{The capacity continuum as a diagnostic tool.}
By treating capacity as an experimental variable rather than a claim, we use the continuum to reveal the \emph{threshold} at which structured divergence emerges.
This threshold has practical implications: if LoRA(r=64) already produces structured divergence, the computational cost of full fine-tuning may not be justified for federated PRM.
If full fine-tuning is required, the finding exposes a fundamental capacity requirement that lighter-weight approaches cannot satisfy.

\textbf{Plan B: If capacity does not affect divergence type.}
If the full capacity continuum shows CD-SPI$_{\text{sym}} \approx 0$ across all configurations (triggering F2), the finding would be equally significant---it would demonstrate that \emph{cross-client PRM divergence is inherently noise-dominated regardless of capacity}, meaning federated PRM aggregation provides no benefit over single-client training.
This negative result would redirect the field toward alternative architectures (MoE, SCAFFOLD) rather than pursuing larger-capacity federated PRMs.

% ============================================================
\section{Limitations}
\label{sec:limitations}
% ============================================================

\begin{enumerate}
    \item \textbf{Architecture scope.} Our causal claims are restricted to dense backbone architectures (Pythia series).
    Extensions to MoE, LLaMA, Qwen, or Mistral architectures require independent validation.
    \item \textbf{Aggregation method.} We only evaluate FedAvg. FedProx, SCAFFOLD, and FedDYN may alter the divergence-to-performance relationship.
    \item \textbf{Heterogeneity regime.} Our domain-based client split (math/code/medical/general) represents one type of Non-IID distribution.
    Dirichlet sampling, label shift, and mixed patterns may produce different CD-SPI profiles.
    \item \textbf{CD-SPI measurement limitations.} CD-SPI operates in cosine space, which is insensitive to embedding magnitude differences.
    The choice of penultimate layer for symmetrical measurement may not capture divergence localized in other layers.
    \item \textbf{Scale.} Experiments use Pythia-1.4B. Scaling laws for CD-SPI structure at 2.8B, 7B, and beyond are unknown.
    \item \textbf{Task specificity.} Our findings are specific to PRM training. Generalization of CD-SPI to other federated learning tasks (e.g., policy model training, image classification) requires independent evaluation.
\end{enumerate}

% ============================================================
\section{Conclusion}
% ============================================================

We introduced CD-SPI, a two-stage diagnostic framework for distinguishing noise-type from structure-type cross-client divergence in federated process reward model training.
Preliminary results from the lowest-capacity configuration (head-only, GELU activation) provide strong confirmation that frozen-backbone PRMs produce zero structured cross-client divergence (CD-SPI$_{\text{sym}} = 0.0011$, EVR $=0.0$, CKA $=1.0$), explaining why FedAvg aggregation under head-only training reduces to averaging identical updates.
The full capacity continuum experiment, currently in progress, will determine whether higher-capacity configurations (LoRA, partial FT, full FT) unlock structured divergence patterns that make federated aggregation meaningful.
The core contribution is the diagnostic framework itself---a signal-structure perspective on federated aggregation that has been overlooked by all existing work on federated reward models.

% ============================================================
\bibliographystyle{plainnat}
\bibliography{refs}
% ============================================================

\appendix
\section{Proofs of Theoretical Properties}
\label{sec:proofs}

\emph{[[TO BE ADDED] --- Formal proofs for Lemmas 1--3 will be provided in the camera-ready version.]}

\section{Full Experimental Configuration Details}

\emph{[[TO BE ADDED] --- Complete YAML configurations, hardware specifications, and reproducibility checklist.]}

\section{Additional Results}

\emph{[[TO BE ADDED] --- Per-round CD-SPI curves, permutation test distributions, CKA matrices, and OOD evaluation results for all configurations.]}

\end{document}
